% !TeX program = pdflatex
% Практичний довідник з курсу «Основи нелінійної динаміки».
% Тема: фрактали.
% Компіляція: pdflatex dovidnyk_practika_fractals.tex (двічі, для змісту).
\documentclass[12pt,a4paper]{article}

\usepackage[T2A]{fontenc}
\usepackage[utf8]{inputenc}
\usepackage[english,ukrainian]{babel}
\usepackage{lmodern}

\usepackage{amsmath,amssymb,amsthm}
\usepackage{mathtools}
\usepackage[a4paper,margin=2.3cm]{geometry}
\usepackage{enumitem}
\usepackage{microtype}
\usepackage{graphicx}
\usepackage{caption}
\usepackage{subcaption}
\usepackage{booktabs}
\usepackage{xcolor}
\usepackage{titlesec}
\usepackage[hidelinks,colorlinks=true,linkcolor=blue!55!black,urlcolor=teal!70!black]{hyperref}
\usepackage{fancyhdr}

% --- Оформлення заголовків ---
\definecolor{accent}{RGB}{30,110,70}
\titleformat{\section}
  {\normalfont\Large\bfseries\color{accent}}{\thesection.}{0.6em}{}
\titleformat{\subsection}
  {\normalfont\large\bfseries\color{accent!85!black}}{\thesubsection.}{0.5em}{}

\pagestyle{fancy}
\fancyhf{}
\renewcommand{\headrulewidth}{0.3pt}
\fancyhead[L]{\small\itshape Практикум: фрактали}
\fancyhead[R]{\small\thepage}

\newcommand{\placeholder}{%
  \par\medskip\noindent\textit{Розв'язок буде додано.}%
}

\addto\captionsukrainian{%
  \renewcommand{\contentsname}{Зміст}%
  \renewcommand{\figurename}{Рисунок}%
}

\graphicspath{{figures/}}
\captionsetup{format=plain,labelfont=bf,justification=centering}
\setlength{\headheight}{14pt}

\begin{document}

% ======================= ТИТУЛЬНА СТОРІНКА =======================
\begin{titlepage}
\thispagestyle{empty}
\centering
\vspace*{1cm}
{\scshape\large Київський національний університет імені Тараса Шевченка\par}
{\scshape факультет комп'ютерних наук та кібернетики\par}
\vspace{1.5cm}
{\color{accent}\rule{\linewidth}{1.2pt}\par}
\vspace{0.6cm}
{\Huge\bfseries Основи нелінійної\\[0.25em] динаміки\par}
\vspace{0.5cm}
{\LARGE Практичний довідник\par}
\vspace{0.35cm}
{\large Фрактали\par}
\vspace{0.4cm}
{\color{accent}\rule{\linewidth}{1.2pt}\par}
\vfill
{\small Київ --- 2026\par}
\end{titlepage}

% ======================= ЗМІСТ =======================
\tableofcontents
\newpage

% =====================================================================
\section{Обов'язкові та додаткові задачі (тема~4)}
% =====================================================================

\subsection{Задача 4.1. Фрактальна розмірність губки Менгера}
Знайдіть фрактальну розмірність губки Менгера.
\placeholder

\subsection{Задача 4.2. Фрактальна розмірність килима Серпінського}
Знайдіть фрактальну розмірність килима Серпінського.
\placeholder

\subsection{Задача 4.3. Збіжність у метриці Хаусдорфа}
Нехай $\{\mathcal{A}_k\}\subset\mathrm{comp}(\mathbb{R}^n)$~--- послідовність компактів,
$\mathcal{A}_k\supset\mathcal{A}_{k+1}$, $k=1,2,\ldots$,
$\mathcal{A}=\bigcap_{k=1}^{\infty}\mathcal{A}_k$. Довести, що
$\mathcal{A}\in\mathrm{comp}(\mathbb{R}^n)$ і в метриці Хаусдорфа
\begin{equation*}
  \lim_{k\to\infty}\mathcal{A}_k=\mathcal{A}.
\end{equation*}
\placeholder

\subsection{Задача 4.4. Відстань Хаусдорфа для множини Кантора}
Нехай $\mathcal{K}$~--- множина Кантора, $\mathcal{K}_0=[0,1]$,
\begin{gather*}
  \mathcal{K}_1=\Bigl[0,\tfrac13\Bigr]\cup\Bigl[\tfrac23,1\Bigr],\\
  \mathcal{K}_2=\Bigl[0,\tfrac{1}{3^2}\Bigr]\cup\Bigl[\tfrac{2}{3^2},\tfrac13\Bigr]
  \cup\Bigl[\tfrac23,\tfrac23+\tfrac{1}{3^2}\Bigr]
  \cup\Bigl[\tfrac23+\tfrac{2}{3^2},1\Bigr],\quad\ldots
\end{gather*}
Знайти відстань Хаусдорфа
$h(\mathcal{K},\mathcal{K}_0)$, $h(\mathcal{K},\mathcal{K}_1)$,
$h(\mathcal{K},\mathcal{K}_2)$, $\ldots$
\placeholder

\subsection{Задача 4.5. Фрактальна розмірність кола}
Знайти фрактальну розмірність множини
$S=\{x\in\mathbb{R}^2:\|x\|=1\}$.
\placeholder

\subsection{Задача 4.6. Фрактальна міра відрізка}
Показати, що фрактальна міра множини $\mathcal{A}=[0,1]$ дорівнює нулю при $d>1$.
\placeholder

\paragraph{Умови для задач 4.7--4.11.}
У $\mathbb{R}^2$ задано одиничний квадрат $\mathcal{K}$ з вершинами в точках
$a^{(1)},a^{(2)},a^{(3)},a^{(4)}$. Позначимо через $F$ систему перетворень
\begin{equation*}
  f_i(x)=r(x-a^{(i)})+a^{(i)},\quad i=1,2,3,4,
\end{equation*}
де $r\in(0,\tfrac12)$, $x\in\mathbb{R}^2$.

\subsection{Задача 4.7. Перші ітерації IFS на квадраті}
Побудувати геометричне представлення перших трьох членів послідовності
$\mathcal{K}_j$, $j=1,2,3$, де $\mathcal{K}_1=\mathcal{K}$,
$\mathcal{K}_{j+1}=f(\mathcal{K}_j)$, $j=1,2,3$,
\begin{equation*}
  f=f_1\cup f_2\cup f_3\cup f_4,\quad
  a^{(1)}=\begin{pmatrix}0\\0\end{pmatrix},\;
  a^{(2)}=\begin{pmatrix}1\\0\end{pmatrix},\;
  a^{(3)}=\begin{pmatrix}0\\1\end{pmatrix},\;
  a^{(4)}=\begin{pmatrix}1\\1\end{pmatrix}.
\end{equation*}
\placeholder

\subsection{Задача 4.8. Алгоритм атрактора IFS}
Розробити алгоритм знаходження атрактора системи функцій $F$.
\placeholder

\subsection{Задача 4.9. Фрактальна розмірність атрактора}
Знайти фрактальну розмірність атрактора, який відповідає системі функцій $F$.
\placeholder

\subsection{Задача 4.10. Розмірність подібності}
Знайти розмірність подібності атрактора системи функцій $F$. Порівняти її
з фрактальною розмірністю.
\placeholder

\subsection{Задача 4.11. Умова відкритої множини}
Перевірити для системи функцій $F$ виконання умов відкритої множини.
Обґрунтувати самоподібність атрактора системи функцій $F$.
\placeholder

\subsection{Задача 4.12. Килим Серпінського}
Розробити алгоритм та побудувати зображення килима Серпінського.
\placeholder

\subsection{Задача 4.13. Крива Коха}
Розробити алгоритм та побудувати зображення кривої Коха.
\placeholder

\subsection{Задача 4.14. П'ятикутник Дарера}
Розробити алгоритм та побудувати зображення п'ятикутника Дарера.
\placeholder

\medskip
\noindent\textit{Додаткові задачі за даною тематикою представлені у [25,~56].}

\newpage
% =====================================================================
\section{Обов'язкові та додаткові задачі (тема~5)}
% =====================================================================

\subsection{Задача 5.1. Період $n\geqslant 3$ і нерухома точка}
Довести, що якщо функція $f$ має орбіту періоду $n\geqslant 3$, то вона має
орбіту періоду~$1$ (тобто має нерухому точку).
\placeholder

\subsection{Задача 5.2*. Періоди $2^n$ та $2^k$}
Довести, що якщо $f$ має орбіту періоду $2^n$, то вона має орбіту періоду
$2^k$ при $n>k$.
\placeholder

\paragraph{Умови для задач 5.3--5.5.}
Нехай $Q(x)=4x(1-x)$, $0\leqslant x\leqslant 1$, та
\begin{equation*}
  T(x)=
  \begin{cases}
    2x, & 0\leqslant x\leqslant \tfrac12,\\[0.3em]
    2(1-x), & \tfrac12<x\leqslant 1.
  \end{cases}
\end{equation*}

\subsection{Задача 5.3*. Топологічна спряженість $Q$ і $T$}
Показати, що відображення $Q$ і $T$ топологічно спряжені на $[0,1]$ за
допомогою $H(x)=\sin^2\!\bigl(\tfrac{\pi}{2}x\bigr)$, тобто
$g=h\circ f\circ h^{-1}$.
\placeholder

\subsection{Задача 5.4. Орбіта періоду 3}
Знайти орбіту $Q(x)$ періоду~$3$, показавши, що
$\bigl\{\tfrac27,\tfrac47,\tfrac67\bigr\}$~--- 3-цикл для $T(x)$.
\placeholder

\subsection{Задача 5.5. Орбіти періодів 4 і 5}
Знайти орбіти $Q(x)$ періодів~4 і~5.
\placeholder

\subsection{Задача 5.7. Орбіта в $\mathbb{C}$}
Для многочлена $f(z)=z^2-1+0.213i$ реалізувати алгоритм знаходження орбіти
на 100 ітерацій при $z_0=0.3+0.2i$ та $z_0=0.31+0.2i$.
\placeholder

\subsection{Задача 5.8. Множина Жюліа ($c=-0.8+0.156i$)}
Реалізувати алгоритм конструювання множини Жюліа для
$f(z)=z^2-0.8+0.156i$.
\placeholder

\subsection{Задача 5.9. Множина Жюліа ($c=0.285+0.01i$)}
Реалізувати алгоритм конструювання множини Жюліа для
$f(z)=z^2+0.285+0.01i$.
\placeholder

\subsection{Задача 5.10. Множина Жюліа ($c=-0.0085+0.71i$)}
Реалізувати алгоритм конструювання множини Жюліа для
$f(z)=z^2-0.0085+0.71i$.
\placeholder

\subsection{Задача 5.11. Множина Мандельброта}
Розробити алгоритм конструювання множини Мандельброта.
\placeholder

\subsection{Задача 5.12. Крива «готель Гільберта»}
Розробити алгоритм та побудувати зображення кривої «готель Гільберта».
\placeholder

\subsection{Задача 5.13. Кардіоїда}
Побудувати зображення кардіоїди за рівнянням
\begin{equation*}
  c=\tfrac12 e^{i\theta}-\tfrac14 e^{2i\theta},\qquad \theta\in[0,2\pi].
\end{equation*}
\placeholder

\subsection{Задача 5.14. Множина Мандельброта для $z^3+c$}
Отримати зображення множини Мандельброта для $f(z)=z^3+c$; показати, що якщо
$|c|>2$, орбіта прямує до $\infty$.
\placeholder

\subsection{Задача 5.15. Басейн $\mathcal{A}(1)$ для $z^4-1$}
Побудувати графічне зображення басейну притягання $\mathcal{A}(1)$ для
$f(z)=z^4-1$.
\placeholder

\subsection{Задача 5.16. Басейн $\mathcal{A}(1)$ для $z^3-z$}
Побудувати графічне зображення басейну притягання $\mathcal{A}(1)$ для
$f(z)=z^3-z$.
\placeholder

\medskip
\noindent\textit{Додаткові задачі за даною тематикою представлені у [25,~60,~62].}

\end{document}
